\section*{THE THEORY OF SETS.}
\phantomsection
\addcontentsline{toc}{section}{THE THEORY OF SETS.}

It can be said that at all times, mathematicians and philosophers have used reasoned arguments from the theory of sets in a more or less conscious way; but in the history of their conceptions of this subject, one must separate sharply all questions linked to the idea of cardinal number (and in particular to the notion of infinity) from those that only introduce the notions of belonging and inclusion. These latter are among the most intuitive and seem never to have raised controversy: it is on them that one can most easily base a theory of syllogism (as Leibniz and Euler were to show us), or axioms such as ``the whole is greater than the sum of its parts'', without talking about that which, in geometry, is concerned with intersections of curves and surfaces. Until the end of the XIXth century, no problem arises in talking about the set (or ``class'' with some authors) of objects possessing such and such a given property; \(^{38}\) and the famous ``definition'' given by Cantor (``By a set is meant a gathering into one whole of objects which are quite distinct in our intuition or our thought'' ({[}47{]}, p.~282)) will give rise, at the time of its publication, to hardly any objections. \(^{39}\) It is altogether different as soon as the notion of set is mixed with that of number or quantity. The question of the infinite divisibility of the line (doubtless posed already by the first Pythagoreans) was, as is known, to lead to considerable philosophical difficulties: from the Eleates to Bolzano and Cantor, mathematicians and philosophers will throw themselves without success against the paradox of the finite quantity made up of an infinity of points without size. It would be of no interest to us to retrace, even summarily, the interminable and impassioned polemics that are aroused by this problem, that constituted a ground particularly favourable to metaphysical or theological digressions; let us note only the point of view at which, from Antiquity, most mathematicians stop. It consists essentially in refusing to debate, short of being able to conclude in an irrefutable way --- an attitude that we will find again among the modern formalists; in the same way as these latter arrange for the elimination of the intervention of ``paradoxical'' sets (see below pp.~31 ff.), the classical mathematicians carefully avoid introducing into their reasoning the ``actual infinity'' (that is to say sets made up of an infinity of objects conceived as existing simultaneously, at least in their thought), and content themselves with ``potential infinity'', that is to say the possibility of increasing every given quantity (or of decreasing it if it is a case of a ``continuous'' quantity). \(^{40}\) If this point of view contained a certain dose of hypocrisy, \(^{41}\) it allowed however the development of the larger part of classical mathematics (including the theory of ratios and later the infinitesimal Calculus); \(^{42}\) it even seemed an excellent safeguard, above all after the quarrels caused by the infinitely small, and had become a dogma almost universally subscribed to until well into the XIXth century.

A first glimmer of the general notion of equipotence appears in a remark of Galileo ({[}122 b{]}, v. VIII, pp.~78-80): he observes that the map \(n \mapsto n^2\) establishes a bijective correspondence between the natural numbers and their squares and that it follows that the axiom ``the whole is greater than the part''

could not be applied to an infinite set. But far from inaugurating a rational study of infinite sets, this remark seems to have had no other effect than to reinforce mistrust of actual infinity; it is already the conclusion of Galileo himself, and Cauchy, in 1833, only quotes him in order to approve of his attitude.

The necessities of Analysis --- and particularly the deep study of functions of real variables, which was taking place during the whole of the XIXth century --- are at the origins of what was to become the modern theory of sets. When Bolzano, in 1817, proves the existence of the lower limit of a set bounded below in R {[}27 c{]}, he still reasons, like the greater part of his contemporaries, ``by comprehension'', speaking not of an arbitrary set of real numbers, but of an arbitrary property of these latter. But when, 30 years later, he writes his ``Paradoxien des Unendlichen'' {[}27 b{]} (published in 1851, three years after his death), he does not hesitate to claim the right to the existence of the ``actual infinity'' and to speak of arbitrary sets. He defines, in this work, the general notion of the equipotence of two sets, and proves that two compact intervals in R are equipotent; he notes also that the characteristic difference between finite sets and infinite sets consists in that an infinite set E is equipotent to a subset distinct from E, but he does not give any convincing proof of this assertion. The general tone of this work is in any case more philosophical than mathematical, and for lack of disassociating sufficiently clearly the notion of the power of a set from that of size or order of infinity, Bolzano fails in his attempts to construct infinite sets of greater and greater power, and lets himself be drawn on this occasion into mixing his reasoning with considerations on divergent series that are totally devoid of sense.

It is due to the genius of G. Cantor that the theory of sets as we know it today was created. He also started with Analysis, and his work on trigonometric series, inspired by those of Riemann, bring him naturally, in 1872, to a first attempt to classify ``exceptional'' sets which occur in this theory, \(^{43}\) by means of the notion of successive ``derived sets'', which he introduces at this point. It is no doubt in connection with these researches, and also with his method for defining the real numbers ({[}47{]}, pp.~92-97), that Cantor begins to be interested in the problems of equipotence, for in 1873, he remarks that the set of rational numbers (or the set of algebraic numbers) is countable; and in his correspondence with Dedekind, which starts about this time {[}48{]}, he is seen posing the question of the equipotence between the set of whole numbers and the set of real numbers, which he succeeds in solving negatively a few weeks later. Then, from 1874 on, it is the problem of dimension that preoccupies him, and for three years he seeks in vain to establish the impossibility of a bijective correspondence between R and \(R^{n}(n > 1)\) , before succeeding, to his great surprise, \(^{44}\) in defining such a correspondence. Having such results as new as they were surprising, he devotes himself entirely to the theory of sets. In a series of 6 memoirs published in the Mathematische Annalen between 1878 and 1884, he tackles both the problems of equipotence, the theory of totally ordered sets, the topological properties of R and of R \(^{n}\) and the problem of measure; and it is admirable to see with what clarity these notions appear in his hands little by little, notions which appear so inextricably tied to the classical conception of the ``continuum''. From 1880, he has the idea of iterating ``transfinitely'' the construction of ``derived sets''; but this idea takes on life only two years later, with the introduction of well ordered sets, one of the most original discoveries of Cantor, thanks to which he can tackle a detailed study of cardinal numbers and formulate the ``problem of the continuum'' {[}47{]}.

It was impossible that such bold conceptions, reversing a twice millennial tradition, and leading to results so unexpected and of such a paradoxical appearance, would be accepted without resistance. In fact, amongst the mathematicians who were influential at that time in Germany, Weierstrass was the only one to follow with some favour the work of Cantor (his erstwhile pupil); this latter was to run up on the other hand against the irreducible opposition of Schwarz and especially of Kronecker. \(^{45}\) It is, it would appear, as much the constant tension generated by the opposition to his ideas, as his fruitless efforts to prove the continuum hypothesis, which caused in Cantor the first symptoms of a nervous illness from which his mathematical productivity would suffer. \(^{46}\) He only really took up again his interest in the theory of sets in 1887, and his last publications date from 1895-1897; he develops there mostly the theory of totally ordered sets and the calculus of ordinals. He had also proved in 1890 the inequality \(m < 2^{m}\) ; however, not only did the problem of the continuum remain without answer, but there remained in the theory of cardinals a more serious gap, for Cantor had been unable to establish the existence of a well ordering between arbitrary cardinals. This gap was going to be filled, on the one hand by the theorem of F. Bernstein (1897) showing that the relations \(a \leq b\) and \(b \leq a\) imply \(a = b\) , \(^{47}\) and above all by the theorem of Zermelo {[}342 a{]} proving the existence of a well order on every set --- a theorem conjectured already in 1883 by Cantor ({[}47{]}, p.~169).

However Dedekind, from the beginning, had not ceased to follow with a sustained interest the researches of Cantor; but while this latter concentrated his attention on infinite sets and their classification, Dedekind followed his own thoughts on the notion of number (which had already led him to his definition of irrationals by ``cuts''). In his opuscule Was sind und was sollen die Zahlen, published in 1888, but the essentials dating from 1872-78 ({[}79{]}, v. III, p.~335), he shows how the notion of natural number (on which, as has been seen, the whole of classical Mathematics had ended up by relying) could itself be derived from fundamental notions in the theory of sets. Developing (no doubt the first explicitly) the elementary properties of arbitrary maps from one set to another (neglected until Cantor, who was only interested in bijective correspondences), he introduces, for every map f of a set E into itself, the notion of ``chain'' of an element \(a \in E\) relative to f, namely the intersection of the sets \(K \subset E\) such that \(a \in K\) and \(f(K) \subset K\) .⁴⁸ He then takes as his definition of an infinite set E the fact that there exists a bijective map \(\phi\) of E into E such that \(\phi(E) \neq E\) ;⁴⁹ if further there exists such a map \(\phi\) and an element \(a \notin \phi(E)\) for which E is the chain of a, Dedekind says that E is ``simply infinite'', remarks that the ``axioms of Peano'' are then satisfied and shows (before Peano) how starting from there all the theorems of elementary arithmetic are obtained. The only missing part of his exposé is the axiom of infinity that Dedekind (following Bolzano) believes himself able to prove by considering the ``world of thought'' of humans (``Gedankenwelt'') as a set.⁵⁰

On another tack, Dedekind had been led by his work on arithmetic (and notably the theory of ideals) to consider the notion of ordered set in a more general setting than Cantor. While this latter restricts himself exclusively to totally ordered sets, \(^{51}\) Dedekind tackles the general case, and makes in particular a deep study of lattices ({[}79{]}, v. II, pp.~236-271). This work was hardly noticed at the time; although these results, found again by several authors, have been the object of numerous publications since 1935, their historical importance arises much less from the possibilities for applications, fairly thin no doubt, of this theory, than from the fact that they constituted one of the first examples of careful axiomatic construction. On the other hand, the first results of Cantor on sets which were countable or had the power of the continuum were rapidly to have numerous and important applications right up to the most classical questions of Analysis \(^{52}\) (without talking naturally of those parts of the Cantorian work that inaugurated General Topology and measure theory; on that see pp.~141 and 221). Moreover, from the last years of the XIXth century the first uses of the principle of transfinite induction were appearing, having become, especially after the proof of the theorem of Zermelo, an indispensable tool in all parts of modern mathematics. Kuratowski was to give, in 1922, an often more amenable version of this principle, avoiding the use of well ordered sets ({[}189 a{]}, p.~89); it is in this form, found again later by Zorn {[}344{]}, that it is mainly used at the present time. \(^{53}\)

Towards the end of the XIXth century, the essential concepts of Cantor had thus gained their cause. \(^{54}\) We have seen that, at about this time, the formalisation of mathematics is achieved and the use of the axiomatic method is almost universally agreed. But, at the same time, a ``crisis of foundations'' of a rare violence was opening up, which was to shake the mathematical world for more than 30 years, and to seem at times to compromise, not only all the recent acquisitions, but even the most classical parts of Mathematics.

\section*{THE PARADOXES OF THE THEORY OF SETS AND THE CRISIS OF FOUNDATIONS.}
\phantomsection
\addcontentsline{toc}{section}{THE PARADOXES OF THE THEORY OF SETS AND THE CRISIS OF FOUNDATIONS.}

The first ``paradoxical'' sets appeared in the theory of cardinals and ordinals. In 1897, Burali-Forti remarked that one cannot consider the existence of a set consisting of all the ordinals, because this set would be well ordered, and so isomorphic to one of its segments distinct from the whole, which is absurd {[}42{]}. \(^{55}\) In 1899, Cantor observes (in a letter to Dedekind) that it is no longer possible to speak of the cardinals forming a set, nor to speak of the ``set of all sets'' without reaching a contradiction (the set of subsets of this latter ``set'' \(\Omega\) would be equipotent to a subset of \(\Omega\) , which is contrary to the inequality \(m < 2^{m}\) ) ({[}47{]}, pp.~444-448). In 1905 finally, Russell analysing the proof of this inequality, shows that the reasoning that establishes it proves (without appealing to the theory of cardinals) that the notion of the ``set of all sets which are not elements of themselves'' is, it also, contradictory {[}266{]}. \(^{56}\)

It could be believed that such ``antinomies'' would only be manifested in the peripheral regions of mathematics, characterised by the consideration of sets of a ``size'' inaccessible to intuition. But other ``paradoxes'' were soon to threaten the most classical parts of mathematics. Berry and Russell ({[}266{]}, v. I, pp.63-64), simplifying an argument of J. Richard {[}258{]}, observe in effect that the set of whole numbers of which the definition can be expressed in less than sixteen words in English is finite, but that it is however contradictory to define a whole number as ``the least whole number which is not definable by less than sixteen words in English'', because this definition consists of fifteen words only.

Although such reasoning, so far from that usually used by mathematicians, may have appeared to many of them as a kind of riddle, they pointed the way none the less to the necessity for a revision of the bases of mathematics, aiming to eliminate ``paradoxes'' of this nature. But although there was unanimity concerning the urgency of this revision, radical divergences were soon to appear concerning the manner of realising it. For a first group of mathematicians, be they ``idealists'', be they ``formalists'', \(^{57}\) the situation created by the ``paradoxes'' of the theory of sets was very analogous to that which resulted, in geometry, from the discovery of non-Euclidean geometries or of ``pathological'' curves (such as curves without tangents); it must lead to a similar, but more general, conclusion namely that it is hopeless to seek to base an arbitrary mathematical theory on an appeal (explicit or not) to ``intuition''. The position can be summed up by the very words of the principal adversary of the formalist school: ``The formalist'' says Brouwer ({[}39 a{]}, p.~83), ``submits that human reason does not have at its disposal exact pictures of straight lines or numbers greater than ten, for example \ldots It is true that certain relations between mathematical entities, that we take as axioms, we can deduce from other relations according to fixed rules, with the conviction that in this way we derive truth from other truths by logical reasoning \ldots{[}But{]} for the formalist, mathematical exactitude resides only in the development of the sequence of relations, and is independent of the meaning that one could wish to give to the relations or entities that they link.''

It consists, therefore, for the formalist, in giving to the theory of sets an axiomatic basis quite analogous to that of elementary geometry, where one does not concern oneself with what the ``objects'' are that are called ``sets'', nor what is meant by the relation \(x \in y\) , but where one enumerates the conditions to be imposed on this latter relation; of course, this must be done in such a way as to include, as far as possible, all the results of the theory of Cantor, all the while making it impossible for ``paradoxical'' sets to exist. The first example of such an axiomatisation was given by Zermelo in 1908 {[}342 c{]}; he avoids sets that are ``too large'' by the introduction of an ``axiom of selection (Aussonderung)'', a property \(P(x)\) only determining a set made up of the elements which have it if \(P(x)\) already entails a relation of the form \(x \in A\) .⁵⁸ But the elimination of paradoxes analogous to ``Richard's paradox'' could only be done by restricting the meaning attached to the notion of ``property''; on that, Zermelo contents himself with describing in very vague terms a type of property that he calls ``defined'', and to indicate that one must limit oneself to these latter in the application of the axiom of selection. This point was made precise by Skolem {[}286 a{]} and Fraenkel {[}113 b{]}; as they observed, its elucidation forces us to place ourselves in a completely formalised system, where the notions of ``property'' and of ``relation'' have lost all ``significance'' and have become simple designations for the collections formed following explicit rules. That necessitates of course that one must put into the system the rules of Logic that are used, which was not yet the case in the system of Zermelo-Fraenkel.

Other axiomatisations of the theory of sets were proposed following on from this. Let us mention principally that of von Neumann {[}342 a and c{]} which, more than the system of Zermelo-Fraenkel, approached the primitive conception of Cantor: this latter, to avoid paradoxical sets, had already proposed ({[}47{]}, pp.~445-448), in his correspondence with Dedekind, to distinguish between two kinds of sets, the ``multiplicities'' (``Vielheiten'') and the ``sets'' (``Mengen'') proper, the second ones being characterised in that they can be thought of as a single object. It is this idea that von Neumann makes precise in distinguishing two types of objects, ``sets'' and ``classes''; in his system (almost completely formalised), classes are distinguished from sets in that they cannot be placed on the left of the sign ∈. One of the advantages of such a system is that it rehabilitates the notion of ``universal class'' used by the logicians of the XIXth century (which naturally is not a set); let us point out also that the system of von Neumann avoids (for the theory of sets) the introduction of axiom schemes, replaced by suitable axioms (which makes its logical study easier). Some variants of the system of von Neumann have been given by Bernays and Gödel {[}130 b{]}.

The elimination of paradoxes seems to have been well carried out by the systems above, but at the cost of restrictions which cannot but seem to be very arbitrary. In favour of the system of Zermelo-Fraenkel, one can say that it restricts itself to formulating restrictions that only sanction the current practice in the applications of the notion of set to the various mathematical theories. The systems of von Neumann and of Gödel are further removed from the usual conceptions; on the other hand, it is not excluded that it is easier to insert certain mathematical theories still in their beginnings in the framework provided by such systems rather than in the narrower framework of the system of Zermelo-Fraenkel.

Certainly one could not affirm that any of these solutions give the impression of being definitive. If they satisfy the formalists, it is that these latter refuse to take into consideration the individual psychological reactions of each mathematician; they assume that a formalised language has completed its task when it can transcribe mathematical reasoning into a form lacking ambiguity, and so useful as a vehicle for mathematical thought; everyone is free, they would say, to think what he likes about the ``nature'' of mathematical objects or the ``truth'' of the theorems that he uses, so long as his reasoning can be transcribed into the common language. \(^{59}\)

In other words, from a philosophical point of view, the attitude of the formalists consists in being disinterested in the problem set by the ``paradoxes'', in abandoning the platonic position which aimed to attribute to mathematical notions an intellectual ``content'' common to all mathematicians. Many mathematicians draw back faced with this break with tradition. Russell, for example, seeks to avoid paradoxes by analysing their structure in a deeper way. Taking up again an idea first brought forward by J. Richard (in the article {[}258{]} where he displayed his ``paradox'') and developed later by H. Poincaré {[}251 c{]}, Russell and Whitehead observe that the definitions of paradoxical sets all violate the following principle, called the ``principle of the vicious circle'': ``An element of which the definition involves the totality of the elements of a set can not belong to that set'' ({[}266{]}, v. I, p.~40). Also it is this statement that serves as basis for the Principia, and it is to satisfy it that the ``theory of types'' is developed in this work. Like that of Frege by which it is inspired, the logic of Russell and Whitehead contains ``propositional variables''; the theory of types proceeds to a classification of these different variables, of which the main lines are the following. Starting from a ``domain of individuals'' which are not specified and which can be qualified as ``objects of order 0'', the relationships where the variables (free or tied) are individuals, are said to be ``objects of the first order''; and in general, the relations in which the variables are objects of order ≤ n (one at least being of order n) are said to be ``objects of order \(n + 1\) ''. \(^{60}\) A set of objects of order n can then only be defined by a relation of order \(n + 1\) , a condition that allows the elimination without trouble of the paradoxical sets. \(^{61}\) But the principle of the ``hierarchy of types'' is so restrictive that by adhering to it strictly, one ends up with a mathematics of inextricable complexity. \(^{62}\) To escape this consequence, Russell and Whitehead are constrained to introduce an ``axiom of reducibility'' affirming the existence, for every relation between ``individuals'', of a relation of the first order to which it is equivalent; a condition quite as arbitrary as the axioms of the formalists, and which reduces considerably interest in the construction of the Principia. So the system of Russell and Whitehead has had more success with the logicians than with the mathematicians; it is in any case not entirely formalised, \(^{63}\) and as a result there are numerous obscure details. Various efforts have been made to simplify and clarify this system (Ramsey, Chwistek, Quine, Rosser); tending to use languages that are more and more fully formalised, these authors replace the rules of the Principia (which still had a certain intuitive foundation) with restrictions only taking account of the wording of the aggregations being considered; not only do these rules appear to be as gratuitous as the prohibitions formulated in the systems of Zermelo-Fraenkel or of von Neumann, but, being further removed from the practice of mathematicians, they have, in many cases, led to unacceptable consequences, which the author had not foreseen (such as the paradox of Burali-Forti, or the negation of the axiom of choice).

For mathematicians of the preceding schools, it was above all a case of not renouncing any part of the heritage of the past: ``From the paradise that Cantor has created for us'', says Hilbert ({[}163 c{]}, p.~274), ``no one should be able to chase us''. To do this, they are willing to accept limitations on mathematical reasoning, hardly essential because consonant with use, but which do not seem to be imposed by our mental customs and the intuitive idea of the notion of set. Anything seems to them preferable to the intrusion of psychology into the criteria of validity in mathematics; and rather than ``putting into account the characteristics of our brains'', as Hadamard says ({[}12{]}, p.~270), they are resigned to imposing on the mathematical domain limits which are for the most part arbitrary, so long as they include classical Mathematics and do not risk hampering future progress.

It is a very different matter when we see the attitude of mathematicians belonging to the tendency of which it remains for us to speak. If the formalists accept the renunciation of the control of the ``eyes of the mind'' so far as mathematical reasoning is concerned, the mathematicians that have been called ``empiricists'', ``realists'' or ``intuitionists'' refuse this abdication; they must have a kind of internal certainty guaranteeing the ``existence'' of the mathematical objects with which they are concerned. As long as it was only a case of renouncing spatial intuition, there was no serious objection, since the arithmetical ``models'' allowed retrenchment behind the intuitive notion of whole numbers. But irreducible opposition manifests itself when it is a matter of reducing the notion of whole number to that (much less intuitively precise) of set, then of imposing on the handling of sets barriers without intuitive foundation. The first in time of these opponents (and the one who, by the authority of his genius, was to exercise the most influence) was H. Poincaré; having admitted, not only the axiomatic point of view for geometry and the arithmetisation of Analysis, but also a good part of the theory of Cantor (which he was one of the first to apply fruitfully in his work), he refuses on the other hand to conceive that arithmetic may be able, it also, to be justified by an axiomatic treatment; the principle of complete induction seems to him in particular a fundamental intuition of our mind, where it is impossible to see a pure convention \(^{64}\) {[}251 c{]}. Hostile in principle to formalised languages, of which he contested the usefulness, he confuses constantly the notion of whole number in formalised mathematics, and the use of whole numbers in the theory of proof, which was starting up then, and of which we will speak later; no doubt it was difficult at that time to make as clearly as today --- after 50 years of study and discussion --- this distinction which however had been clearly felt by a Hilbert or a Russell.

The critics of this kind multiplied after the introduction of the axiom of choice by Zermelo in 1904 {[}342 a{]}. Its use in many previous proofs in Analysis or the theory of sets had until then gone by almost unnoticed; \(^{65}\) it is in following an idea suggested by Erhard Schmidt that Zermelo, explicitly formulating this axiom, deduced from it in an ingenious way a satisfactory proof of the existence of a well ordering on every set. Coming at the same time as the ``paradoxes'', it appears that this new way of reasoning, by its odd allure, had thrown many a mathematician into confusion; it is enough to see the strange misunderstandings which surged up about this topic, in the following volume of Mathematische Annalen, under the signature of mathematicians who were as familiar with the Cantorian methods as Schoenflies and F. Bernstein. The criticisms of E. Borel, published in this same volume, are more substantial and attach themselves clearly to the point of view expressed by H. Poincaré about whole numbers; they are developed and discussed in an exchange of letters between E. Borel, Baire, Hadamard and Lebesgue, who had remained in the classical tradition in French mathematics {[}12{]}. Borel starts by denying the validity of the axiom of choice because it contains in general an uncountable infinity of choices, which is inconceivable intuitively. But Hadamard and Lebesgue observe that a countable infinity of arbitrary successive choices is no more intuitive, since it consists of an infinity of operations, that it is impossible to conceive as really taking place. For Lebesgue, who enlarged the debate, everything reduces to knowing what is meant by saying that a mathematical object ``exists''; it is necessary, for him, that one ``names'' explicitly a property defining it in a unique way (a ``functional'' property, we would say); for a function such as the one used by Zermelo in his reasoning, that is what Lebesgue calls a ``law'' of choice; if, he continues, this requirement is not satisfied, and that ones limits oneself to ``thinking'' about this function instead of ``naming'' it, can one be sure that during the reasoning process one is always thinking of the same thing ({[}12{]}, p.~267)? In any case, this brings Lebesgue to new doubts; already the question of the choice of a single element in a set seems to him to raise difficulties: one must be sure that such an element ``exists'', that is to say that one can ``name'' at least one of the elements of the set. \(^{66}\) Can one then speak of the ``existence'' of a set of which one does not know how to ``name'' each element? Already Baire does not hesitate to deny the ``existence'' of the set of all subsets of a given infinite set (loc. cit., pp.~263-264); in vain does Hadamard observe that these requirements lead to renouncing even the possibility of speaking of the real numbers: it is to just that conclusion that E. Borel finally agrees. Putting aside the fact that countability seems to have acquired freehold, one has returned just about to the classical position of the adversaries of the ``actual infinite''.

All these objections were not very systematic; it was down to Brouwer and his school to undertake a complete remould of mathematics guided by similar, but even more radical, principles. We would not dare to summarise here a doctrine as complex as intuitionism, which consists as much of psychology as mathematics, and we will limit ourselves to indicating some of the most striking aspects, referring for more details to the work of Brouwer himself {[}39 a and b{]} and to the exposition by Heyting {[}162{]}. For Brouwer, Mathematics is identical with the ``exact'' part of our thought, based on the first intuition of the sequence of natural numbers, and which it is impossible to translate without mutilation into a formal system. It is in any case only ``exact'' in the thought of mathematicians, and it is a chimera to hope to develop an instrument of communication between them which is not subject to all the imperfections and ambiguities of language; at most it is possible to hope to awaken in the interlocutor a favourable state of mind by descriptions that are more or less vague ({[}162{]}, pp.~11-13). Intuitionist mathematics attaches hardly more importance to logic than to language: a proof is not conclusive by virtue of logical rules fixed once and for all, but because of the ``immediate evidence'' of each of its steps. This ``evidence'' must moreover be interpreted in a more restrictive way than by E. Borel and his followers: it is thus that in intuitionist mathematics, it cannot be said that a relation of the form ``R or (not R)'' is true (principle of the excluded middle), unless for all system of values given to the variables appearing in R, one can prove that one of the two propositions R, ``not R'' is true; for example, from the equation ab = 0 between two real numbers, it cannot be concluded that ``a = 0 or b = 0'', for it is easy to create explicit examples of real numbers a, b for which one has ab = 0 without knowing, at the actual time, how to prove either of the two propositions a = 0, b = 0 ({[}162{]}, p.~21).

It is not astonishing that, starting from such principles, the intuitionist mathematicians finish up with results that are very different from the classical theorems. A whole part of these disappear, for example most of the ``existential'' theorems in Analysis (such as the theorems of Bolzano and Weierstrass for real functions); if a function of a real variable ``exists'' in the intuitionist sense, it is ipso facto continuous; a bounded monotone sequence of real numbers does not necessarily have a limit. On the other hand, many classical notions become ramified, for the intuitionist, into several distinct fundamental notions: there are thus two notions of convergence (for a sequence of real numbers) and eight of countability. It goes without saying that transfinite induction and its applications to modern Analysis are (as for the greater part of the theory of Cantor) condemned without appeal.

It is only in this way that, according to Brouwer, mathematical propositions can acquire a ``content''; the formalist reasonings that go beyond what the intuitionists allow are judged to be without value, since one can no longer give them a ``meaning'' to which the intuitive notion of ``truth'' could be applied. It is clear that such judgements can only rely on a pre-existing notion of ``truth'', of a psychological or metaphysical nature. It is to say in practice that they escape all discussion.

It is undoubted that the vigorous attacks that came from the intuitionist camp have sometimes forced not only the mathematical schools in the avant-garde, but even the followers of traditional mathematics to go on the defensive. A famous mathematician has admitted having been impressed by the attacks to the point of having voluntarily restricted his work to branches of mathematics judged to be ``certain''. But such cases can hardly have been frequent. The intuitionist school, of which the memory is no doubt destined to remain only as an historical curiosity, would at least have been of service by having forced its adversaries, that is to say definitely the immense majority of mathematicians, to make their position precise and to take more clearly notice of the reasons (the ones of a logical kind, the others of a sentimental kind) for their confidence in mathematics.

\section*{METAMATHEMATICS}
\phantomsection
\addcontentsline{toc}{section}{METAMATHEMATICS}

The absence of contradiction has at all times been considered as a sine qua non condition for all mathematics, and from the time of Aristotle, logic was sufficiently developed so that it was perfectly well known that from an inconsistent theory anything could be deduced. The proofs of ``existence'', considered as indispensable since Antiquity, visibly had no other aim than to guarantee that the introduction of a new concept did not risk entailing a contradiction, particularly when this concept was too complicated to fall immediately under ``intuition''. We have seen how this requirement became more pressing with the arrival of the axiomatic point of view in the XIXth century and how the construction of arithmetical ``models'' responded to it. But arithmetic itself, could it be inconsistent? A question that without doubt would not have been thought of before the end of the XIXth century, so much did the whole numbers appear to belong to what was most certain in our intuition; but after the ``paradoxes'' everything seemed to have been questioned, and one can understand that the feeling of insecurity that they created had brought mathematicians, around 1900, to bear with more care on the problem of the non-contradiction of arithmetic, so as to save at least classical mathematics from shipwreck. So this problem is the second of those that Hilbert enumerated in his famous lecture at the international Congress of 1900 ({[}163 a{]}, v. III, pp.~229-301). In doing this, he laid down a new principle that was to have great repercussions: whereas in traditional logic the non-contradiction of a concept only made it ``possible'', it was equivalent for Hilbert (at least for mathematical concepts defined axiomatically) to the existence of the concept. This implied apparently the necessity of proving a priori the non-contradiction of a mathematical theory before even being able to develop it legitimately; it is just in this way that H. Poincaré understands it when, to make a breach in formalism, he takes to his account the idea of Hilbert, underlining with a mischievous pleasure the extent to which the formalists were far, at that stage, from being able to achieve it ({[}251 c{]}, p.~163). We will see later how Hilbert was to pick up the gauntlet; but before that, it must be noted here that, under the influence of his authority and that of Poincaré, the requirements set down by this latter were to be for a long time accepted without reservation as well by the formalists as by their adversaries. One consequence was the belief, very widespread among the formalists, that the theory of proof of Hilbert was an integral part of mathematics, of which it constituted the indispensable prolegomena. This dogma does not seem justified to us, \(^{67}\) and we consider that the intervention of metamathematics in the exposition of logic and mathematics can and should be reduced to the very elementary part that is concerned with the handling of abbreviating symbols and criteria of deduction. It is thus not a case, contrary to what Poincaré thinks, of ``claiming the freedom from contradiction'', but rather to consider, with Hadamard, that the absence of contradiction, even when it is not proven, is noted ({[}12{]}, p.~270).

It remains for us to give a brief historical sketch of the efforts of Hilbert and his school, and to retrace summarily, not only the evolution that was finally to lead to the negative result of Gödel and justify a posteriori the scepticism of Hadamard, but also all the progress that followed on from this, concerning the knowledge of the mechanism of mathematical reasoning, and which makes modern metamathematics an autonomous science of incontestable interest.

Already in 1904, in a lecture at the international Congress ({[}163 c{]}, pp.~247-261), Hilbert attacks the problem of consistency in arithmetic. He observes first that there is no question of proving it by making use of a model, \(^{68}\) and outlines the principles of another method: he proposes considering the true propositions of formalised arithmetic as collections of signs without meaning, and to prove that in using the rules governing the formation and linkage of these collections, it is impossible to obtain a collection which is a true proposition and whose negation is also a true proposition. He sketches even a proof of this nature for a formalism less extensive than that of arithmetic; but, as was observed by H. Poincaré soon after ({[}251 c{]}, p.~185), this proof uses essentially the principle of recurrence, and so appears to rely on a vicious circle. Hilbert did not answer this criticism immediately, and about fifteen years passed by without anyone attempting to develop his ideas; it is in 1917 only that (motivated by the wish to answer the attacks of the intuitionists) he harnesses himself again to the problem of the foundations of mathematics, which he will not cease from now on to consider until the end of his scientific career. In his work on this subject, which is spread out from 1920 to 1930 approximately, and in which a whole school of young mathematicians take part (Ackermann, Bernays, Herbrand, von Neumann), Hilbert brings out little by little the principles of his ``theory of proof'' in a more precise way: acknowledging implicitly the validity of the criticism of Poincaré, he admits that in metamathematics, the arithmetical reasoning that is used can only be based on our intuition of the whole numbers (and not on formalised arithmetic); to do this it appears to be essential to restrict this reasoning to ``finite procedures'' (``finite Prozesse'') of a type allowed by the intuitionists: for example, a proof by contradiction cannot prove the metamathematical existence of a collection or sequence of collections, there must be an explicit law of construction. \(^{69}\) On the other hand Hilbert enlarges in two directions his initial programme: not only does he tackle the consistency of arithmetic, but he aspires also to prove the consistency of the theory of real numbers and even that of the theory of sets; \(^{70}\) further, to the problems of consistency come to be added those of independence of axioms, of completeness and of decision. We will pass quickly in review these various questions and point out the principal research to which they have given rise.

The proof of the independence of a system of propositions \(A_{1}, A_{2}, \ldots, A_{n}\) consists in showing that, for each index i, \(A_{i}\) is not a theorem in the theory \(T_{i}\) obtained by taking as axioms the \(A_{j}\) with index \(j \neq i\) . This point will be established if one knows a consistent theory \(T_{i}^{\prime}\) in which the \(A_{j}(j \neq i)\) are theorems, as well as ``not \(A_{i}\) ''; thus one can consider the problem in two ways depending on whether or not one allows that certain theories (such as arithmetic or the theory of sets) are consistent. In the second case, one has to deal with a problem of ``absolute'' consistency. On the contrary, the first type of problem reduces to a problem of ``relative'' consistency, by means of a construction of appropriate ``models'', and many proofs of this kind have been thought up, well before mathematics had taken up a completely formalised aspect: it is enough to remember the models of non-Euclidean geometry (see p.~133), the questions of the independence of the axioms of elementary geometry dealt with by Hilbert in the ``Grundlagen der Geometrie'' {[}163 c{]} as well as the work of Steinitz on the axiomatisation of Algebra and those of Hausdorff and his successors on that of Topology.

A theory T is said to be complete if, for every proposition A of T (not containing any letter other that the constants of T), one of the propositions A, (not A), is a theorem of \(T.^{71}\) . If one puts on one side certain very rudimentary formalisms, of which the completeness is easily proved ({[}164{]}, p.~35), the results obtained in this area are essentially negative; the first in time, is due to K. Gödel {[}130 a{]} who showed that, if T is consistent and if the axioms of formalised arithmetic are theorems of T, then T is not complete. The fundamental idea of his ingenious method consists in establishing a bijective correspondence (of course, by means of ``finite processes'') between the metamathematical statements and certain propositions in formalised arithmetic; we will restrict ourselves to a broad sketch. \(^{72}\) For each collection A which is a term or a relation of T, one begins by associating (by a procedure which is an explicit construction, which can be applied quasi mechanically) a whole number \(g(A)\) , in a bijective way. In the same way, to each proof D of T (considered as a succession of collections) one can associate a whole number \(h(D)\) in a bijective way. Finally, one can give a procedure for constructing explicitly a relation \(P(x,y,z)\) of T, \(^{73}\) such that, in T, \(P(x,y,z)\) requires that x,y,z are whole numbers, and satisfies the following two conditions:

\(1^{\circ}\) if \(D\) is a proof of \(A(\lambda)\) , where \(A(x)\) is a relation of \(\mathcal{T}\) , and \(\lambda\) is an explicit whole number (that is to say a term of \(\mathcal{T}\) which is a whole number), then \(P(\lambda, g(A(x)), h(D))\) is a theorem of \(\mathcal{T}\) ;

\(2^{\circ}\) if the explicit whole number \(\mu\) is not of the form \(h(D)\) , or if \(\mu = h(D)\) and if \(D\) is not a proof of \(A(\lambda)\) , then (not \(P(\lambda, g(A(x)), \mu)\) ) is a theorem of \(\mathcal{T}\) .

Then let \(S(x)\) be the relation (not \((\exists z)P(x,y,z)\) ), and let \(\gamma = g(S(x))\) , which is a term of \(\mathcal{T}\) . If \(\mathcal{T}\) is consistent, there is no proof of the proposition \(S(\gamma)\) in \(\mathcal{T}\) . In fact, if \(D\) were such a proof, \(P(\gamma,g(S(x)),h(D))\) would be a theorem of \(\mathcal{T}\) ; but this relation is none other than \(P(\gamma,\gamma,h(D))\) , and it follows that \((\exists z)(P(\gamma,\gamma,z)\) would also be a theorem of \(\mathcal{T}\) ; as this last relation is equivalent to (not \(S(\gamma)\) ), \(\mathcal{T}\) would be inconsistent. On the other hand what has just been said shows that for every explicit whole number \(\mu\) , (not \(P(\gamma,\gamma,\mu)\) ) is a theorem of \(\mathcal{T}\) . The result is that there is no proof, in \(\mathcal{T}\) , of (not \(S(\gamma)\) ), for this last relation is equivalent to \((\exists z)P(\gamma,\gamma,z)\) and the existence of a whole number \(\mu\) such that \(P(\gamma,\gamma,\mu)\) would require that \(\mathcal{T}\) were inconsistent, by virtue of what comes before. \(^{74}\) This metamathematical theorem of Gödel has been subsequently generalised in various directions ({[}181{]}, chap.~XI). \(^{75}\)

The relation \(S(\gamma)\) of T of which it is shown thus that there is no proof in T either of this relation, or of its negation, is obviously made up for the needs of the cause and does not tie up in any natural way with any mathematical problem. Much more interesting is the fact that if \(\mathcal{T}\) denotes the theory of sets (with the system of axioms of von Neumann-Bernays), neither the continuum hypothesis nor its negation are provable in \(\mathcal{T}\) . This remarkable result was established in two stages: in 1940, Gödel proved that the theory obtained by adjoining to \(\mathcal{T}\) the continuum hypothesis \(2^{\aleph_0} = \aleph_1\) was not inconsistent {[}130 b{]}; then, in 1963, P. Cohen proved the same when one adjoins to \(\mathcal{T}\) the relation \(2^{\aleph_0} = \aleph_2\) (or \(2^{\aleph_0} = \aleph_n\) for an arbitrary whole number \(n > 1\) ) {[}67{]}.

The decision problem (``Entscheidungsproblem'') is without doubt the most ambitious of all those which are stated in metamathematics: it consists in knowing whether, for a given formal language, one can imagine a ``universal procedure'' which is quasi mechanical which, when applied to any relation of the formalism in question, indicates in a finite number of operations whether that relation is true or not. \(^{76}\) The solution of this problem was already part of the grand design of Leibniz, and it appears that at one time the school of Hilbert believed its realisation was quite close. It is a fact that one can describe such procedures for formalisms containing few primitive symbols and axioms ({[}181{]}, pp.~136-141). But the efforts made to make precise the decision problem, by outlining exactly what is meant by ``universal procedure'' have until now resulted only in negative results ({[}181{]}, pp.432-439). Besides, the solution of the decision problem for a theory T entrains immediately the knowledge of whether T is consistent or not, since it suffices to apply the ``universal procedure'' to a relation of T and its negation; and we will see that it is impossible to solve the question in this way for the usual mathematical theories. \(^{77}\)

It is indeed in the question of the consistency of mathematical theories --- the origin and very heart of Metamathematics --- that the results have revealed themselves to be the most deceiving. During the years 1920-1930, Hilbert and his school had developed new methods for attacking these problems; after having proved the consistency of partial formalisms, covering part of arithmetic (cf.~{[}158{]}, {[}324 d{]}), they believed they had reached the goal and proved, not only the consistency of arithmetic, but also that of the theory of sets, when Gödel, relying on the incompleteness of arithmetic, deduced from it the impossibility of proving, by the ``finite procedures'' of Hilbert, the consistency of every theory T containing this latter. \(^{78}\)

The theorem of Gödel does not however shut the door completely on attempts to prove consistency, so long as one abandons (at least partially) the restrictions of Hilbert concerning ``finite procedures''. It is thus that Gentzen in 1936 {[}127{]}, succeeds in proving the consistency of formalised arithmetic, by using ``intuitively'' transfinite induction up to the countable ordinal \(\varepsilon_{0}\) .⁷⁹ The value of the ``certainty'' that one can attach to such reasoning is without doubt less convincing than for that which satisfies the initial requirements of Hilbert, and is essentially a matter of personal psychology for each mathematician; it remains no less true that similar ``proofs'' using ``intuitive'' transfinite induction up to a given ordinal, would be considered important progress if they could be applied, for example, to the theory of real numbers or to a substantial part of the theory of sets.

\chapter{NOTATION; COMBINATORIAL ANALYSIS.}

History and archaeology make us conscious of a large number of ``systems of notation''; their initial goal was to attach to each individual whole number (up to a limit which depends on the needs of the application) a name and a written representation, made up of a combination of a restricted number of symbols, accomplished by following more or less regular laws. By far the most frequent procedure consists in splitting up the whole numbers into sums of ``successive units'' \(b_{1}, b_{2}, \ldots, b_{n}, \ldots\) , of which each is a multiple by a whole number of the preceding one; and if in general \(b_{n}/b_{n-1}\) is taken to be equal to the same number b (the ``base'' of the system, most often 10), many exceptions to this rule are noticed, as with the Babylonians, where \(b_{n}/b_{n-1}\) is sometimes equal to 10, sometimes to 6 {[}232{]}, and in the chronological system of the Mayas, where \(b_{n}/b_{n-1}\) is equal to 20 except for n = 2, and where \(b_{2}/b_{1} = 18\) {[}228{]}. As for the corresponding written representation, it must indicate the number of ``units'' \(b_{i}\) of each order i; in many systems (such as with the Egyptians, the Greeks and the Romans), successive multiples \(k.b_{i}\) , (where k varies from 1 to \((b_{i+1}/b_{i}) - 1\) ) are represented by symbols which depend at the same time on k and on i. A first and important improvement consisted in representing all the numbers \(k.b_{i}\) (for the same value of k) by the same symbol: it is the principle of the ``notation of position'', where the index i is indicated by the fact that the symbol representing \(k.b_{i}\) appears in the ith position in the sequence of ``slices'' making up the number being represented. The first system of this kind is found among the Babylonians, who, without doubt already 2000 years before Christ, denote by the same sign all the multiples \(k.60^{\pm i}\) corresponding to arbitrary values of the exponent i ({[}232{]}, pp.~93-109). The inconvenience of such a system resides of course in the ambiguity of the symbols used, so long as nothing indicates that the units of a certain order may be absent, in other words, so long as the system is not completed by the introduction of a ``zero''. One sees however the Babylonians making do without such a symbol for the greater part of their history; in fact they only use such a ``zero'' in the last two centuries before Christ, and still only in the interior of a number; until then, the context alone had to specify the significance of the number under consideration. Only two other systems use a ``zero'' systematically: that of the Mayas (in use, it would appear, already from the beginning of the Christian era {[}228{]}), and our actual decimal system, which (by the agency of the Arabs) is derived from Hindu mathematics, where its use is attested already from the first centuries of our era. It must be noted moreover that the conception of zero as a number (and not as a simple symbol of separation) and its introduction into calculations, also count amongst the original contributions of the Hindus ({[}78{]}, v. I). Of course, once the principle of the notation of position has been acquired, it was easy to extend it to any base whatever; the discussion concerning the merits of different ``bases'' proposed since the XVIIth century is a matter of the techniques of numerical Calculation and should not be tackled here. We limit ourselves to remarking that the operation that serves as foundation to these systems, the so-called ``Euclidean'' division, does not appear before the Greeks, and no doubt goes back to the first Pythagoreans, who made it the essential tool of their theoretical Arithmetic (see p.~85).

The general problems of enumeration, grouped under the name of ``combinatorial Analysis'', do not seem to have been tackled before the last centuries of classical Antiquity: only the formula \(\binom{n}{2}=n(n-1)/2\) is attested to in the IIIrd century of our era. The Hindu mathematician Bhaskara (XIIth century) knows the general formula for \(\binom{n}{p}\) . A more systematic study is to be found in a manuscript of Levi ben Gerson, at the beginning of the XIIIth century: he obtains the recurrence formula allowing the calculation of the number \(V_{n}^{p}\) of arrangements of n objects p at a time, and in particular the number of permutations of n objects; he states also rules equivalent to the relations \(\binom{n}{p}=V_{n}^{p}/p!\) and \(\binom{n}{n-p}=\binom{n}{p}\) ({[}311{]}, v. VI, pp.~64-65). But this manuscript seems to have been ignored by his contemporaries, and its results have only been rediscovered little by little by the mathematicians of the following centuries. Among the later advances, let us pick out that Cardan proves that the number of non-empty subsets of a set of n elements is \(2^{n}-1\) ; Pascal and Fermat, setting up the calculus of probability, rediscover the expression for \(\binom{n}{p}\) , and Pascal is the first to observe the relationship between these numbers and the binomial formula: this latter seems to have been already known to the Arabs from the XIIIth century, the Chinese in the XIVth century, and it was rediscovered in the West at the beginning of the XVIth century, as well as the method of calculation by recurrence known as the ``arithmetical triangle'', that is normally attributed to Pascal ({[}311{]}, v. VI, pp.~35-38). Finally Leibniz, around 1676, obtains (without publishing it) the general formula for ``multinomial coefficients'', rediscovered independently and published 20 years later by de Moivre.

\chapter{THE EVOLUTION OF ALGEBRA.}

There are few notions, in Mathematics, that are more primitive than that of a law of composition: it seems inseparable from the first rudiments of calculation with whole numbers and measurable quantities. The most ancient documents that remain from the Mathematics of the Egyptians and Babylonians show us that they are already in possession of a complete system of rules of calculation for whole numbers \textgreater{} 0, rational numbers \textgreater{} 0, lengths and areas; even though the Babylonian texts that have reached us only deal with problems in which the data have explicit numerical values, \(^{1}\) they leave no doubt as to the generality attributed to the rules that are used, and denote a quite remarkable technical facility in the handling of equations of the first and second degrees ({[}232{]}, pp.~179 and following). No hint of a justification of the rules used is found in any case, nor even a precise definition of the operations that occur: these like those remain in the domain of empiricism.

A similar concern, on the other hand, shows itself already very clearly among the Greeks of the classical era; one does not meet yet, it is true, an axiomatic treatment of the theory of whole numbers (such an axiomatisation will not appear until the end of the XIXth century; see p.~24); but, already in the Elements of Euclid, many are the passages giving formal proofs of rules of calculation quite as intuitively ``obvious'' as those of calculation with whole numbers (for example, the commutativity of the product of two rational numbers). The most remarkable proofs of this kind are those that are connected with the theory of quantities, the most original creation of Greek Mathematics (equivalent, as we know, to our theory of real numbers \textgreater{} 0; see pp.~148 ff.): Euclid considers there, amongst others, the product of two ratios of quantities, shows that it is independent of the form in which the two ratios are given (first example of a ``quotient'' of a law of composition by an equivalence relation), and that it is commutative ({[}107{]}, Book V, prop. 22-23). \(^{2}\)

We must not hide, though, the fact that this progress towards rigour is accompanied, in Euclid, with a stagnation, and even some points of regression, where the techniques of algebraic calculations are concerned. The crushing preponderance of Geometry (in the context of which the theory of quantities is manifestly conceived) paralyses all autonomous development of an algebraic notation: the elements occurring in the calculations must, at every moment, be ``represented'' geometrically: over and above, the two laws of composition which occur are not defined on the same set (addition of ratios is not defined in a general way, and the product of two lengths is not a length but an area); there is, as a result, a lack of flexibility which makes it almost impractical to handle algebraic relations of degree greater than the second.

It is only with the decline of classical Greek Mathematics that one sees Diophantus return to the traditions of the ``logistitians'' or professional calculators, who had continued to apply as they stood the rules inherited from the Egyptians and Babylonians: not burdening themselves any more with geometrical representations for those ``numbers'' that he considers, he is naturally led to develop rules for abstract algebraic calculation; for example, he gives rules which (in modern language) are equivalent to the formula \(x^{m+n} = x^{m}x^{n}\) for small values (positive or negative) of m and n ({[}91 a{]}, v. I, pp.~8-13); a bit further on (pp.12-13) is found the statement of the ``rule of signs'', first hint of calculations with negative numbers; finally Diophantus uses, for the first time, a literal symbol to represent an unknown in an equation. On the other hand, he seems barely bothered to connect to general ideas the methods which he applies to the solving of his problems; as to the axiomatic conception of laws of composition, such as was beginning to dawn with Euclid, it appears foreign to the thought of Diophantus as well as to his immediate followers; it will not be found again in Algebra before the beginning of the XIXth century.

It was necessary first, during the intermediary centuries, that on the one hand a system of algebraic notation was developed, adequate for the expression of abstract laws, and that, on the other hand, the notion of ``number'' received a broadening that allowed, by the observation of sufficiently diversified special cases, an ascent to general concepts. On that point, the axiomatic theory of ratios of quantities, created by the Greeks, was insufficient, because it only made precise the intuitive notion of a real number \textgreater{} 0 and the operations on these numbers, already known to the Babylonians in a more confused form; it will be a case on the contrary, now, of ``numbers'' of which the Greeks had not conceived, and of which, at first, no sensible ``representation'' presented itself: on the one hand, zero and negative numbers, which appeared from the high Middle Ages in Hindu Mathematics; on the other, imaginary numbers, creation of the Italian algebraists in the XVIth century.

If one puts aside the zero, introduced first of all as a symbol of notation before being considered as a number (see p.~45), the common character of these extensions is to be (at least at first) purely ``formal''. By that it must be understood that the new ``numbers'' appear first of all as the results of operations applied in situations where they do not have, according to their strict definition, any meaning (for example, the difference a - b of two whole numbers when a \textless{} b): whence the name of ``false'', ``fictitious'', ``absurd'', ``impossible'', ``imaginary'', etc. numbers, that have been attributed to them. For the Greeks of the classical era, taken up above all with clear thinking, such extensions were inconceivable; they could only come from calculators more disposed than were the Greeks to have faith even though somewhat mystical in the power of their methods (``the generality of Analysis'', as would say the XVIIIth century), and to allow themselves to be dragged along by the mechanics of their calculations without checking at each step its soundness; a faith justified besides most often a posteriori, by the exact results to which the extensions, to these new mathematical objects, of rules of calculation valid only, in all rigour, for numbers known previously, led. That explains how there is a gradually increasing boldness in considering for themselves (independently of all applications to concrete calculations), these generalisations of the notion of number, which, at first, only arose as intermediaries in a sequence of operations of which the point of departure and the goal were true ``numbers''; once this step was taken, more or less tangible interpretations of these new entities began to be sought, which acquired thus the right to be quoted in Mathematics. \(^{4}\)

On this point, the Hindus are already conscious of the interpretation that negative numbers must have in certain cases (a debt in a commercial problem, for instance). In the following centuries, as there is a diffusion into the West (by the intermediary of the Arabs) of the methods and results of Greek and Hindu mathematics, one becomes more used to the handling of these numbers, and one begins to have other ``representations'' for them which are geometric or dynamic. There in any case, together with a progressive improvement in algebraic notation, is the only notable progress in Algebra during the end of the Middle Ages.

At the beginning of the XVIth century, Algebra took new flight, thanks to the discovery, by the mathematicians of the Italian school, of the solution ``by radicals'' of the equation of the 3rd, then of that of the 4th degree (see pp.~72 ff.); it is on this occasion that, in spite of their repugnance, they found themselves so to speak constrained to introduce into their calculations imaginary numbers; little by little, however, confidence in calculation with these ``impossible'' numbers is born, as with negative numbers, and that although no ``representation'' of them was thought up for more than two centuries.

On the other hand, algebraic notation received its decisive improvements from Viète and from Descartes; starting with this latter, algebraic notation is already, with a few exceptions, that which we use today.

From the middle of the XVIIth century to the end of the XVIIIth century, it appears that the vast horizons opened up by the creation of the infinitesimal Calculus made to some extent for a neglect of Algebra in general, and in particular mathematical thought on the laws of composition or on the nature of real and complex numbers. \(^{5}\) It is thus that the composition of forces and the composition of velocities, well known in Mechanics already from the end of the XVIIth century, did not exert any repercussion on Algebra, although they contained already in embryo vector calculus. Indeed we must await the movement of ideas which, around 1800, leads to the geometrical representation of complex numbers (see p.~161), to see, in pure Mathematics, the use of vector addition. \(^{6}\)

It is at about the same time that, for the first time in Algebra, the notion of a law of composition extends. in two different directions, to elements that no longer come associated with ``numbers'' (in the widest meaning given up to that point to this word) except by distant analogies. The first of these extensions is due to C. F. Gauss, at the time of his arithmetical researches on quadratic forms \(ax^{2} + bxy + cy^{2}\) with integral coefficients. Lagrange had defined, in the set of forms with the same discriminant, an equivalence relation, \(^{7}\) and had, on the other hand, proved an identity which supplied, in this set, a commutative law of composition (not defined everywhere); starting from these results, Gauss shows that this law is compatible with a slightly different equivalence relation ({[}124 a{]}, v. I, p.~272): ``One sees by that'', he then says, ``what should be understood by a class composed from two or more classes''. He proceeds afterwards to the study of the ``quotient'' law that he has thus just defined, establishes in substance that it is (in modern language) an abelian group law, and this by arguments of which the generality surpasses by far, in the majority of cases, the special case that Gauss considers (for example, the argument by which he proves the uniqueness of the symmetric element can be applied to any law of composition whatever (ibid., p.~273)). But he does not stop there: returning a little later to the question, he recognises the analogy between the composition of classes and the multiplication of integers modulo a prime number \(^{8}\) (ibid. p.~371), but notes also that the group of classes of quadratic forms of given discriminant is not always a cyclic group; the indications that he gives on this matter make us think that he might perhaps have realised, at least in this particular case, the general structure of finite abelian groups ({[}124 a{]}, v. I, p.~374 and v. II, p.~266).

The other series of researches of which we wish to speak ends, it also, with the notion of group, thus introducing it in a definitive way into Mathematics: it is the ``theory of substitutions'', a development of the ideas of Lagrange, Vandermonde and Gauss on the solution of algebraic equations. We will not give here a detailed history of this question (see pp.~75 ff.); we must remember the definition by Ruffini {[}265{]}, then Cauchy ({[}56 a{]}, (2), v. I, p.~64), of the ``product'' of two permutations of a finite set, \(^{9}\) and the first notions about finite groups of permutations: transitivity, primitivity, neutral element, commuting elements, etc.. But these first researches remain, on the whole, fairly superficial, and it is Evariste Galois who must be considered the real initiator of the theory: having, in his memorable work {[}123{]}, reduced the study of algebraic equations to that of the groups of permutations that he associates with them, he deepens considerably this latter, as much in its concern with general properties of groups (it is Galois who first defines the notion of normal subgroup and recognises its importance), as in the determination of groups possessing particular properties (where the results that he obtains are classed still today amongst the most subtle of the theory). It is also to

Galois that the first idea of the ``linear representation of groups'' is due, \(^{10}\) and this fact proves clearly that he was in possession of the notion of the isomorphism of two group structures, independently of their ``realisations''.

However, if it appears incontestable that the brilliant methods of Gauss and Galois had led them to a very wide conception of the law of composition, they did not have the opportunity to develop particularly their ideas on this point, and their work did not have any immediate action on the evolution of abstract Algebra. \(^{11}\) It is in a third direction that the clearest progress towards abstraction is made: following reflections on the nature of imaginaries (of which the geometrical representation had caused, at the beginning of the XIXth century, fairly substantial work), the algebraists of the English school bring out first, from 1830 to 1850, the abstract notion of a law of composition, and enlarge immediately the field of Algebra by applying this notion to a host of new mathematical objects: the algebra of Logic with Boole (see p.~8), vectors, quaternions and general hypercomplex systems with Hamilton {[}145 a{]}, matrices and non-associative laws with Cayley ({[}58{]}, v. I, pp.~127 and 301, and v. II, pp.~185 and 475). A parallel evolution takes place independently on the continent, notably as far as vector Calculus is concerned (Möbius, Bellavitis), linear Algebra and hypercomplex systems (Grassmann) (see pp.~61). \(^{12}\)

From all this bubbling up of original and fruitful ideas that come to revive Algebra in the first half of the XIXth century, this latter emerges renewed even in its trends. Before, methods and results gravitated around the central problem of the solution of algebraic equations (or diophantine equations in Number Theory): ``Algebra'', says Serret in the Introduction to his Higher Algebra Course {[}284{]}, ``is, properly speaking, the Analysis of equations''. After 1850, even if the treatises on Algebra still leave for a long time the preeminence to the theory of equations, new research is no longer dominated by the concern for immediate applications to the solution of numerical equations, and is more and more oriented towards what we would now consider as the essential problem of Algebra, the study of algebraic structures as such.

This work is divided fairly neatly into three streams, which extend respectively the three movements of ideas that we have analysed above, and which are pursued in parallel without mutual measurable influence until the last years of the XIXth century. \(^{13}\)

It is first the building up, by the German school of the XIXth century (Dirichlet, Kummer, Kronecker, Dedekind, Hilbert) of the theory of algebraic numbers, coming out of the work of Gauss, to whom is due the first study of this kind, that of numbers \(a + bi\) ( \(a\) and \(b\) rational). We will not follow here the evolution of this theory: we must simply pick up, for our purposes, the abstract algebraic notions that emerge there. Already from the first successors of Gauss, the idea of a field (of algebraic numbers) is at the base of all the work in question (as are also the researches of Abel and Galois on algebraic equations); its field of application grows when Dedekind and Weber {[}80{]} copy the theory of algebraic functions of one variable onto that of algebraic numbers. Also due to Dedekind ({[}79{]}, v. III, p.~251) is the introduction of the notion of ideal, which provides a new example of a law of composition between sets of elements; it is to him and Kronecker that the ever greater role played by abelian groups and modules in the theory of algebraic fields goes back; we will return to this later on (see pp.~63 ff.~and 93 ff.).

We leave also till later the history of the development of linear Algebra (p.~57) and hypercomplex systems (p.~117), which take place without introducing any new algebraic notion during the end of the XIXth century and the beginning of the XXth century, in England (Sylvester, W. Clifford) and in America (B. and C. S. Peirce, Dickson, Wedderburn) following the route laid out by Hamilton and Cayley, in Germany (Weierstrass, Dedekind, Frobenius, Molien) and in France (Laguerre, E. Cartan) in a manner independent of the Anglo-Saxons and following fairly different methods.

As for the theory of groups, it is mainly under the form of the theory of finite permutation groups that it is developed at first, following the publication of the works of Galois and their diffusion by the works of Serret {[}284{]}, and above all the great ``Treatise on Substitutions'' of C. Jordan {[}174 a{]}. This latter summarises there, perfecting them greatly, the work of his predecessors on the particular properties of permutation groups (transitivity, primitivity, etc.), obtaining results of which the greater part have hardly been surpassed since; he studies as well, in a deep way, very important particular groups, the linear groups and their subgroups (see p.~138); as well, it is he who introduces the fundamental notion of the representation of one group by another, as well as (a little later) that of quotient group, and who proves part of the theorem known by the name of ``Jordan-Hölder Theorem''. \(^{14}\) It is finally to Jordan that the first study of infinite groups goes back {[}174 b{]}, that S. Lie on the one hand, F. Klein and H. Poincaré on the other, were to develop considerably, in two different directions, a few years later.

In the meantime, it had gradually become obvious that what is essential in a group is its law of composition and not the nature of the objects making up the group (see for example {[}58{]}, v. II, pp.~123 and 131 and {[}79{]}, v. III, p.~439). However, even the researches on finite abstract groups are still for a long time conceived as studies about groups of permutations, and it is only around 1880 that there is a conscious development of the autonomous theory of finite groups. We cannot follow further the history of this theory; we limit ourselves to mentioning two of the tools which still today are amongst the most used in the study of finite groups, and which both go back to the XIXth century: the theorems of Sylow \(^{15}\) on p-groups, which date from 1872 {[}303{]} and the theory of characters, created in the last years of the century by Frobenius ({[}119{]}, v. III, pp.1-37). We send the reader wishing to deepen his knowledge of the theory of finite groups and the numerous and difficult problems it gives rise to, to the monographs {[}44 b{]}, {[}131{]}, {[}291{]} and {[}341{]}.

It is towards 1880 also that presentations of groups begin to be studied systematically; before that, only presentations of particular groups had been studied, for example of the alternating group \(\mathfrak{A}_5\) in the work of Hamilton {[}145 b{]}, or of the monodromy groups of differential equations and Riemann surfaces (Schwarz, Klein, Schläfli). W. Dyck was the first {[}97{]} to define (without giving it a name) the free group generated by a finite number of generators, but he is less interested in it for itself than as a ``universal'' object allowing the precise definition of what is meant by a group given ``by generators and relations''. Developing an idea first put out by Cayley, and visibly influenced elsewhere by the work of the school of Riemann mentioned above, Dyck describes an interpretation of a group with a given presentation, where each generator is represented by a product of two inversions with respect to tangent circles or secants (see also {[}44 b{]}, chap.~XVIII). A little later, after the development of the theory of automorphic functions by Poincaré and his successors, then the introduction by Poincaré of the tools of algebraic Topology, the first studies of fundamental groups will go hand in hand with that of group presentations (Dehn, Tietze), the two theories providing mutual support. It is in any case a topologist, J. Nielsen who, in 1924 introduced the terminology free group in the first deep study of its properties {[}234{]}; almost immediately afterwards E. Artin (still in connection with topological questions) introduces the notion of free product of groups, and O. Schreier defines more generally the free product with amalgamated subgroup. Here again, we cannot tackle the history of subsequent developments in this direction, referring to the work {[}216{]} for more details.

It is not either the place to speak of the extraordinary prosperity that has been known, since the end of the XIXth century, by the idea of group (and that of invariant, to which it is intimately linked) in Analysis, in Geometry, in Mechanics and in theoretical Physics. It is by an analogous invasion of this notion, and of the algebraic notions that are related to it (groups with operators, rings, ideals, modules) in those parts of Algebra which until then appeared fairly remote from their proper domain, that the last period of the evolution that we are retracing here is marked, and which ends with the synthesis of the three tendencies that we have followed above. This unification is above all the work of the modern German school: begun with Dedekind and Hilbert in the last years of the XIXth century, the work of axiomitisation of Algebra was vigorously pursued by E. Steinitz {[}294 a{]}, then, from 1920, under the impulsion of E. Artin, E. Noether and the algebraists of their schools (Hasse, Krull, O. Schreier, van der Waerden). The treatise of van der Waerden {[}317 a{]}, published in 1930, gathered together for the first time these works in a complete exposition, opening the way and serving as guide to the multiple subsequent researches into abstract Algebra.
